% =============================================
% CHAPITRE 3 : CONCEPTION
% =============================================
\chapter{Conception}

\section{Introduction}
L'\'{e}tape de conception traduit le cahier des charges en sp\'{e}cifications techniques
et architecturales. Elle sert de plan au d\'{e}veloppement. Ce chapitre pr\'{e}sente
les choix architecturaux et la mod\'{e}lisation de la base de donn\'{e}es.

\section{Architecture Globale du Syst\`{e}me}

L'application \textbf{PT\_BACK} est d\'{e}velopp\'{e}e sous forme d'une \textbf{API RESTful} structur\'{e}e
sur le design pattern \textbf{MVC (Mod\`{e}le-Vue-Contr\^{o}leur)}, adapt\'{e} pour r\'{e}pondre
\`{a} des requ\^{e}tes \texttt{HTTP} et servir le format JSON pluto\^{o}t que des vues rendues.

\vspace{0.4cm}
\begin{itemize}[leftmargin=*]
    \item \textbf{Client} : Les applications Front-end ou Mobiles (ex: React, VueJs).
    \item \textbf{Auth Server (Keycloak)} : G\`{e}re avec ind\'{e}pendance les comptes
    et \'{e}met les \textit{JSON Web Tokens (JWT)}.
    \item \textbf{Backend (PT\_BACK/Laravel)} : Intercepte les requ\^{e}tes, v\'{e}rifie
    la l\'{e}gitimit\'{e} du JWT et traite la logique m\'{e}tier des projets.
\end{itemize}

\section{Diagramme de Cas d'Utilisation}

Le diagramme de cas d'utilisation illustre de mani\`{e}re synth\'{e}tique les
interactions des acteurs avec les principaux modules de l'API.

\begin{itemize}[leftmargin=*, label=\textcolor{emsiRed}{\ding{118}}]
    \item \textbf{Manager} $\rightarrow$ Cr\'{e}ation/Gestion des projets ; Import/Export de donn\'{e}es.
    \item \textbf{Chef de Projet} $\rightarrow$ Suivi des projets ; Cr\'{e}ation et validation des t\^{a}ches.
    \item \textbf{D\'{e}veloppeur} $\rightarrow$ Consultation des t\^{a}ches assign\'{e}es ; Mise \`{a} jour des statuts.
\end{itemize}

\section{Mod\'{e}lisation des Donn\'{e}es (Base de donn\'{e}es relationnelle)}

L'impl\'{e}mentation r\'{e}elle sous Laravel utilise les mod\`{e}les \textit{Eloquent ORM}.
Nous avons extrait la sch\'{e}matique \`{a} partir des migrations effectu\'{e}es dans le framework.

\subsection{Entit\'{e}s Principales}
\begin{itemize}[leftmargin=*]
    \item \textbf{Projects} : \texttt{id}, \texttt{name}, \texttt{client}, \texttt{start\_date}, \texttt{status}, \texttt{manager\_id}, \texttt{chef\_de\_projet\_id}.
    \item \textbf{Tasks} : \texttt{id}, \texttt{title}, \texttt{goal}, \texttt{status}, \texttt{project\_id}.
    \item \textbf{Managers}, \textbf{ChefDeProjets}, \textbf{Developers} : Repr\'{e}sentent les diff\'{e}rentes
    entit\'{e}s humaines au sein du syst\`{e}me.
\end{itemize}

\subsection{Associations et Tables de Jonction (Pivots)}
Plusieurs relations "*Many-to-Many*" s'op\`{e}rent dans l'application :
\begin{itemize}[leftmargin=*]
    \item \textbf{developer\_project} : Associe les D\'{e}veloppeurs aux Projets. Elle embarque
    les champs suppl\'{e}mentaires \texttt{position} et \texttt{joined\_at}.
    \item \textbf{developer\_task} : Associe les D\'{e}veloppeurs aux T\^{a}ches en permettant
    la sp\'{e}cification d'un \texttt{role} et le timestamp \texttt{assigned\_at}.
    \item \textbf{sla\_projects} et \textbf{sla\_tasks} : Entit\'{e}s permettant de logguer et
    lister les SLAs d\'{e}sir\'{e}s.
\end{itemize}

\section{Conclusion}
La rigueur accord\'{e}e \`{a} la mod\'{e}lisation garantira que l'impl\'{e}mentation dans
Laravel s'effectuera de mani\`{e}re fluide, en s'appuyant sur ses relations \textit{Eloquent}.
Ceci nous m\`{e}nera droit aux r\'{e}alisations effectives.
